% critique.tex — f167v, vv. 423–465 : commentaire critique
% Articles textuels tirés du commentaire de l'auteur (latex/Aristophane-Thesmo Harmattan 2023 *.tex),
% texte et notes conservés ; répartition détaillée dans repartition.md.
\cc{433–42}{}{Ce chant du chœur trouvera son écho aux vers 520-530 : semblable réplique au point de vue métrique\nc{À une exception près : au vers 437b où le texte est probablement corrompu.} qui permettait aux deux passages d'être chantés sur la même mélodie.}
\cc{446}{ἐμοὶ γὰρ ἁνὴρ ἀπέθανεν μὲν ἐν Κύπρῳ.}{La précision du lieu de décès du mari (ἁνὴρ ἀπέθανεν ἐν Κύπρῳ) n'a pas manqué d'intriguer les modernes scholiastes.
Leur inclination naturelle étant portée au scabreux, ils ont pour une bonne partie voulu y voir un tour signifiant que la matrone athénienne, mère de cinq enfants, avait tué son mari au déduit avec ses appétits démesurés ; «~mourir de sa belle mort~», comme l'heureux \italique{Félix Faurenicator}. 
L'hypothèse est plaisante mais gratuite et donc à rejeter ! tous les lexiques consultés ne soufflent mot d'une telle expression qu'on aurait nécessairement relevée. 
D'autres se sont essayés à prendre ce vers littéralement. Là encore, ils se sont heurtés à un mur : les dernières expéditions sur l'île de Chypre feraient de notre procurette une septuagénaire au moins !.. et comme il ne semble pas y avoir de jeux de mots, il ne reste qu'une seule voie : une erreur de transcription du texte. 
(C'est d'autant plus vraisemblable que l'expression « ἀπέθανεν ἐν Κύπρῳ », si elle s'était véritablement trouvée dans le texte original, n'aurait pas manqué d'intriguer les commentateurs anciens qui, même sans y rien connaître, auraient ouvert les vannes de leur verve intarissable... au contraire : ce vers n'a étrangement pas laissé la moindre scholie ancienne derrière lui.)\\
L'intervention de la tresseuse de couronnes n'a donc rien de burlesque : elle essaie au contraire d'introduire du pathétique. Or, qu'y a-t-il de plus banal pour apitoyer son auditoire à peu de frais que de sortir le lieu-commun de la femme qui se retrouve veuve alors qu'elle est enceinte (ἔγκυος) ? 
\citComm{ἐμοὶ γὰρ ἁνὴρ ἀπέθανεν μὲν \textbf{ἔγκυος}\\
Mon mari mourut alors que j'étais enceinte \\
Me laissant cinq enfants à charge etc. etc.}}
\cc{463}{πολύπλοκον.}{Les diverses corrections qui ont été proposées pour faire rentrer de force ce vers dans le mètre trochaïque ne me paraissent pas indispensables : on peut simplement considérer qu'il n'y a pas de \italique{correptio attica} et scander :\citComm{καὶ πολύπ|λοκον νόημ᾽, οὐδ᾽\\
$ - \cup - \cup - \cup - - $}}
